\documentclass[11pt,a4paper]{article}
\usepackage[T1]{fontenc}
\usepackage[utf8]{inputenc}
\usepackage[english]{babel}
\usepackage{lmodern,microtype}
\usepackage[a4paper,margin=2.25cm]{geometry}
\usepackage{amsmath,amssymb,booktabs,tabularx,longtable,array}
\usepackage{xcolor,hyperref,enumitem,fancyhdr,titlesec,tcolorbox}
\hypersetup{colorlinks=true,linkcolor=blue!40!black,citecolor=blue!40!black,urlcolor=blue!40!black}
\definecolor{deepblue}{HTML}{183B56}
\definecolor{softblue}{HTML}{EAF2F8}
\definecolor{softamber}{HTML}{FFF4D6}
\definecolor{softred}{HTML}{FBE9E7}
\definecolor{softgreen}{HTML}{E8F5E9}
\pagestyle{fancy}\fancyhf{}\lhead{ED-POMDP Research Edition v1.4}\rhead{Milestone 2R}\cfoot{\thepage}
\setlength{\parindent}{0pt}\setlength{\parskip}{0.55em}\setlength{\emergencystretch}{2em}
\setlist{leftmargin=1.5em,itemsep=0.25em,topsep=0.25em}
\titleformat{\section}{\large\bfseries\color{deepblue}}{\thesection}{0.65em}{}
\titleformat{\subsection}{\normalsize\bfseries\color{deepblue}}{\thesubsection}{0.65em}{}
\newcommand{\go}{\textsc{go}}
\newcommand{\cgo}{\textsc{conditional go}}
\newcommand{\nogo}{\textsc{no-go}}

\title{\textbf{Milestone 2R: Theory and Claim Reset after a Negative Matched-Budget Benchmark}\\
\large Decision-relevant evidence acquisition for software release assurance under partial observability}
\author{Guillaume Harbonnier\\Independent industrial research programme / STRAT-Q}
\date{2 August 2026}

\begin{document}
\maketitle

\begin{abstract}
The Step 2 frozen synthetic benchmark of Evidence-Driven Partially Observable Markov Decision Processes (ED-POMDP) did not support the original broad claim that decision-aware evidence acquisition generally reduces terminal release-decision loss at matched budget. It also did not support the broad claim that explicit evidence-quality modelling generally improves final decisions, although one narrow expected-calibration-error signal was retained. The benchmark evaluated seven policies over 3,360 episodes, generated by four evidence regimes, four fixed budgets, thirty untouched seeds, and seven policies. A deterministic mechanism analysis then formed 2,400 ED-POMDP-versus-baseline pairs, obtained from the same four regimes, four budgets, thirty seeds, and five confirmatory baselines. Of 1,119 pairs in which ED-POMDP improved Brier score, 1,054 retained the same terminal action. The central result is therefore not merely statistical non-significance but a structural conversion failure: better probabilistic beliefs were usually absorbed by unchanged decision regions. Milestone 2R turns that falsification result into a bounded theory revision. It preserves the underlying problem---joint uncertainty about system state and evidence-production quality---while withdrawing broad superiority claims, auditing the complete belief-to-action architecture, and defining conditional hypotheses, progression gates, anti-leakage rules, and a redesigned experimental sequence. This paper explains why the reset is a principled consequence of the preregistered programme rather than a post-hoc displacement of its objectives.
\end{abstract}

\begin{tcolorbox}[colback=softblue,colframe=deepblue,boxrule=0.8pt,arc=2mm]
\textbf{Stage-gate conclusion.} Step 2 closes the original broad superiority claims as unsupported within the frozen benchmark. The programme may continue only through a theory-and-claim reset. It must not continue by tuning the same simulator, increasing the number of seeds, or reinterpreting calibration gains as decision gains.
\end{tcolorbox}

\section{Scientific status after Step 2}

Step 2 evaluated two principal propositions:
\begin{enumerate}
  \item decision-aware value of information should reduce terminal decision loss relative to simpler acquisition policies at the same evidence budget;
  \item explicit modelling of evidence-production quality should improve calibration and, in turn, final release decisions under evidence degradation and model misspecification.
\end{enumerate}

Neither broad proposition met its promotion criterion. No terminal-decision-loss contrast survived Holm correction. Across the four acquisition baselines directly associated with the value-of-information claim, aggregate terminal-loss directions were 10 favourable, 16 adverse, and 38 equal. One degraded-evidence, budget-two ECE contrast survived Holm correction, but its paired bootstrap interval crossed zero and posterior support was sparse. The correct dispositions are therefore:
\begin{itemize}
  \item \texttt{CLM-VOI-001}: \texttt{NOT\_SUPPORTED\_STEP2};
  \item \texttt{CLM-EQ-001}: \texttt{BROAD\_FORM\_NOT\_SUPPORTED\_NARROW\_ECE\_SIGNAL}.
\end{itemize}

These labels are bounded to the frozen simulator, horizons, thresholds, evidence structure, and loss model. They do not assert universal impossibility. They do prohibit a general superiority narrative.

\section{Reconstructing the benchmark counts}

The raw Step 2 benchmark contained
\[
4\ \text{regimes}\times 4\ \text{budgets}\times 30\ \text{seeds}\times 7\ \text{policies}=3{,}360
\]
episode rows.

The Step 2.8 mechanism analysis compared ED-POMDP with five frozen confirmatory baselines for every regime--budget--seed cell:
\[
4\times 4\times 30\times 5=2{,}400
\]
paired diagnostics.

This distinction matters. The 3,360 count describes all executed policy episodes. The 2,400 count describes pairwise mechanism records with ED-POMDP as the reference policy. Among those pairs, ED-POMDP improved Brier score in 1,119 cases, and 1,054 of those improvements retained the same terminal action:
\[
\frac{1{,}054}{1{,}119}=94.19\%.
\]

The risk-only comparison was especially revealing: the two policies selected the same terminal action in all 480 regime--budget--seed pairs, despite frequently different acquisition traces and posterior values.

\section{What Step 2 established}

The result is more informative than ``no significant difference.'' The study identified a break in the value chain:
\[
\text{evidence acquisition}
\rightarrow\text{observation}
\rightarrow\text{belief update}
\not\Rightarrow\text{different and better action}.
\]

A probabilistic improvement is decision-relevant only when it can:
\begin{enumerate}
  \item arise from a credible and sufficiently discriminating observation;
  \item move the belief by a material amount;
  \item cross a decision boundary or alter the stopping decision;
  \item enable a different admissible terminal action or compensating control;
  \item reduce realised or expected loss after evidence cost and delay.
\end{enumerate}

Step 2 frequently achieved the first two conditions without achieving the last three. Better calibration was therefore often operationally inert.

\section{What remains valid}

\subsection{The problem statement remains valid}
Release assurance still combines uncertainty about the system and uncertainty about the evidence-production process. A high pass rate can be misleading when data, environment, oracle, instrumentation, or provenance are weak. Step 2 did not refute that distinction.

\subsection{The frozen architecture is not validated}
The particular Step 2 coupling among acquisition policy, belief update, fixed horizon, unit evidence cost, terminal thresholds, loss ratios, and three-way release action did not produce the expected general advantage.

\subsection{Safety evidence remains exploratory}
ED-POMDP produced zero unsafe \go{} decisions in the three structurally avoidable regimes, compared with two per regime for the classical POMDP. This pattern may be important for a Test Authority, but it was outside the confirmatory p-value family and cannot be promoted retrospectively.

\section{Why Milestone 2R is not goal displacement}

A negative benchmark can create a legitimate concern: is the programme revising theory, or merely moving the target after an unfavourable result? Four safeguards distinguish the present reset from post-hoc rationalisation.

\begin{enumerate}
  \item \textbf{Falsification was part of the original programme.} Step 2 was explicitly designed to expose the claims to matched-budget refutation rather than to demonstrate them by construction.
  \item \textbf{The failed claims remain visible and unpromoted.} The reset does not rename or erase the adverse result; the canonical registry retains the original IDs and dispositions.
  \item \textbf{The new hypotheses are conditional and narrower.} They are not restatements of the old claim with more favourable wording. Each identifies a specific mechanism and a corresponding failure condition.
  \item \textbf{New confirmation requires new data-generating randomness.} Step 2 headline seeds remain immutable historical evidence and may not become a tuning set.
\end{enumerate}

The revised objective is therefore not ``prove ED-POMDP superior somehow.'' It is:
\begin{quote}
Under which conditions does additional or better-qualified evidence change a governed engineering decision, and does that changed decision reduce a materially important loss at comparable total cost?
\end{quote}

\section{Revised register of scientific bets}

\begin{longtable}{p{0.10\textwidth}p{0.38\textwidth}p{0.22\textwidth}p{0.22\textwidth}}
\toprule
ID & Bet & Status after Step 2 & Decision \\
\midrule
P0 & System risk and evidence quality should be represented separately. & Not invalidated; no operational-gain proof. & Retain as a modelling hypothesis. \\
P1 & VOI generally improves decisions at equal budget. & Not supported. & Withdraw the broad form. \\
P2 & Evidence-quality modelling generally improves calibration and decisions. & Broad form not supported. & Decompose into narrower claims. \\
P3 & Better calibration naturally converts into better action. & Contradicted by the observed mechanism. & Withdraw as an implicit assumption. \\
P4 & Benefit depends on the joint coupling of belief, thresholds, loss, stopping, control, and action. & New mechanism hypothesis. & Test explicitly. \\
P5 & Evidence-quality awareness may reduce dangerous release actions. & Descriptive signal only. & Pre-register a dedicated safety endpoint. \\
\bottomrule
\end{longtable}

\section{Milestone 2R work programme}

\subsection{2R.1 Reconstruct the causal chain}

The full decision chain is decomposed as
\[
\text{acquisition}\rightarrow\text{observation}\rightarrow\text{belief}
\rightarrow\text{boundary or stopping}\rightarrow\text{action}\rightarrow\text{loss}.
\]
For every link, Milestone 2R must define necessary conditions, failure modes, diagnostics, and counterexamples.

\subsection{2R.2 Audit structural assumptions}

The reset must examine:
\begin{itemize}
  \item derivation and sensitivity of terminal thresholds;
  \item safety, delay, commercial, and evidence-cost ratios;
  \item fixed horizons versus adaptive stopping;
  \item equal total budget versus cost-sensitive acquisition;
  \item unit costs versus heterogeneous evidence costs;
  \item posterior boundary reachability;
  \item practical indistinguishability among policies;
  \item correlated and redundant evidence;
  \item direct use of evidence quality in terminal admissibility;
  \item explicit semantics and effectiveness of compensating controls for \cgo{}.
\end{itemize}

\subsection{2R.3 Candidate conditional hypotheses}

\textbf{H3-A --- Joint terminal decision.} A terminal rule using both $P(S=\mathrm{bad})$ and evidence-quality belief reduces weighted loss under degraded evidence relative to a risk-only rule.

\textbf{H3-B --- Adaptive stopping.} An adaptive policy reduces mean evidence cost without increasing terminal loss or unsafe-\go{} rate.

\textbf{H3-C --- Compensated conditional release.} A \cgo{} action linked to explicit, verifiable compensating controls reduces loss relative to an unstructured intermediate decision region.

\textbf{H3-D --- Conditional value of information.} VOI provides benefit only when available evidence is heterogeneous, sufficiently discriminating, non-redundant, and capable of crossing a relevant decision boundary.

\textbf{H3-E --- Safety.} Evidence-quality-aware terminal rules reduce unsafe-\go{} rate under controlled degradation at comparable total evidence cost.

These are candidate hypotheses, not promoted claims. Their exact estimands, multiplicity family, effect thresholds, and refutation criteria must be frozen before confirmation.

\subsection{2R.4 H3-D falsifiability gate}

H3-D is not confirmatory until its four qualifiers are operationalised under a frozen protocol. Provenance heterogeneity is computed from a preregistered evidence-family taxonomy and dependency graph. Discriminative power uses Jensen--Shannon divergence with natural logarithms: the raw value is retained in nats and normalized by $\ln 2$ to $D\in[0,1]$. Conditional non-redundancy uses $I(S;O_a\mid O_H)$ in nats, normalized by residual state entropy when that entropy exceeds a frozen $\epsilon_H$, giving $U\in[0,1]$. Boundary reachability is the predictive probability $B\in[0,1]$ that acquisition changes the admissible terminal action or the continue/stop decision.

A cell is eligible only if the provenance rule holds and $D\geq\delta_D$, $U\geq\delta_U$, and $B\geq\delta_B$. The thresholds may be calibrated only on new Milestone 3A development seeds and must be frozen before any Milestone 3B confirmatory seed is opened. The primary estimand is an eligibility-by-policy interaction. Numerical minimum counts are required for both eligible and ineligible groups; failure of either minimum yields \texttt{NOT\_TESTABLE\_IN\_BENCHMARK}, not retrospective reclassification. The complete gate is recorded in \texttt{benchmark/protocol/H3D\_OPERATIONALISATION\_GATE.md}.

\section{Epistemic taxonomy: NONE to SYNTHETIC}

The transition from evidence level \texttt{NONE} to \texttt{SYNTHETIC} records the \emph{kind of adjudicating evidence}, not positive support. Three fields must remain independent:
\begin{itemize}
  \item \textbf{Evidence level}: source class, such as none, formal, synthetic, industrial, or operational;
  \item \textbf{Evidence polarity}: supportive, mixed, adverse, or untested;
  \item \textbf{Disposition}: active, narrowed, unsupported, blocked, or deferred.
\end{itemize}

After Step 2, it would be incorrect to retain \texttt{NONE}: a frozen synthetic benchmark now exists and has adjudicated the claims. It would also be incorrect to interpret \texttt{SYNTHETIC} as support. For \texttt{CLM-VOI-001}, the correct combination is synthetic evidence, adverse/mixed polarity, and \texttt{NOT\_SUPPORTED\_STEP2}. This separation is a central anti-drift rule.

\section{Revised experimental sequence}

\begin{longtable}{p{0.15\textwidth}p{0.74\textwidth}}
\toprule
Milestone & Purpose and progression gate \\
\midrule
2R & Theory and claim reset. No new performance claim. Exit requires audited assumptions, formal candidate mechanisms, and a frozen design boundary. \\
3A & Exploratory mechanism benchmark using new development seeds only. No confirmatory claim. \\
3B & New preregistered benchmark with untouched confirmatory seeds and a small primary endpoint family: weighted terminal loss, unsafe GO, and total evidence cost. \\
4 & Realistic STRAT-Q scenarios with heterogeneous costs, correlated tests and OTEL evidence, compensating controls, governed provenance, schedule constraints, and project-specific loss. \\
5 & Governed retrospective replay and shadow evaluation, gated by data readiness and prior bounded validation. \\
\bottomrule
\end{longtable}

Brier score and ECE remain useful mechanism diagnostics. They are secondary endpoints and may not substitute for decision value.

\section{Threats to validity}

The Step 2 result remains synthetic and benchmark-specific. The state space is small, likelihoods are designed rather than industrially estimated, evidence channels are simplified, fixed horizons suppress optimal stopping, and one loss structure governs all scenarios. The post-hoc mechanism analysis is deterministic and reproducible but descriptive. Conversely, the reset itself risks overfitting the theory to one negative benchmark. The safeguards are narrow conditional hypotheses, explicit non-claims, new development seeds, untouched confirmatory seeds, independent review, and a requirement that simple governed baselines remain first-class comparators.

\section{Conclusion}

Step 2 did not destroy the ED-POMDP research programme. It invalidated an overly simple promise: that selecting and qualifying evidence more intelligently would naturally produce better decisions. The serious scientific output is a mechanism finding: probabilistic improvement often fails to cross the belief-to-action boundary.

\begin{tcolorbox}[colback=softgreen,colframe=green!45!black,boxrule=0.8pt,arc=2mm]
\textbf{Revised central hypothesis.} Decision value does not arise from belief quality alone. It arises from the complete governed architecture that converts belief into stopping, compensating control, and terminal action under asymmetric loss, heterogeneous evidence cost, and hard constraints.
\end{tcolorbox}

Continuation is justified, but only through Milestone 2R. Continuing the former roadmap without this reset would discard the principal scientific value of Step 2.

\begin{thebibliography}{99}
\bibitem{kaelbling1998planning} Kaelbling, L. P., Littman, M. L., and Cassandra, A. R. Planning and acting in partially observable stochastic domains. \emph{Artificial Intelligence}, 101(1--2), 99--134, 1998.
\bibitem{howard1966information} Howard, R. A. Information value theory. \emph{IEEE Transactions on Systems Science and Cybernetics}, 2(1), 22--26, 1966.
\bibitem{fenton2012risk} Fenton, N. and Neil, M. \emph{Risk Assessment and Decision Analysis with Bayesian Networks}. CRC Press, 2012.
\bibitem{lauri2023robotics} Lauri, M., Hsu, D., and Pajarinen, J. Partially observable Markov decision processes in robotics: A survey. \emph{IEEE Transactions on Robotics}, 39(1), 21--40, 2023.
\end{thebibliography}

\end{document}
